\documentclass[11pt,a4paper]{article}
\usepackage[T1]{fontenc}
\usepackage{lmodern,amsmath,amssymb,amsthm,mathtools,booktabs,microtype}
\usepackage[margin=26mm]{geometry}
\usepackage[colorlinks=true,linkcolor=blue,citecolor=blue,urlcolor=blue]{hyperref}
\setlength{\emergencystretch}{2em}
\newtheorem{theorem}{Theorem}[section]
\newtheorem{proposition}[theorem]{Proposition}
\newtheorem{lemma}[theorem]{Lemma}
\newtheorem{corollary}[theorem]{Corollary}
\theoremstyle{definition}\newtheorem{definition}[theorem]{Definition}
\theoremstyle{remark}\newtheorem{remark}[theorem]{Remark}
\newcommand{\Q}{\mathbb Q}\newcommand{\Z}{\mathbb Z}
\newcommand{\R}{\mathbb R}\newcommand{\A}{\mathbb A}
\newcommand{\PP}{\mathbb P}\newcommand{\sqcl}[1]{[#1]}
\DeclareMathOperator{\Tr}{Tr}\DeclareMathOperator{\Nm}{N}
\DeclareMathOperator{\N}{N}\DeclareMathOperator{\rad}{rad}
\DeclareMathOperator{\sfpart}{sf}
\title{Norm and Trace Graphs Have No Quantifier-Free Definitions}
\author{}\date{16 November 2026}
\begin{document}\maketitle
\begin{abstract}
For a proper finite extension $L/K$ of characteristic-zero fields,
the norm and trace graphs, viewed as subsets of $L^2$ with one
$L$-valued input, are Zariski dense. Neither graph admits a
quantifier-free definition in the ring language with coefficients
in $L$. In particular, neither map is a polynomial function in
one variable over $L$. The assertion concerns quantifier-free
definitions; it does not rule out existential definitions with
auxiliary variables.
\end{abstract}

\section{Graphs in the ambient extension field}
Let $L/K$ be a proper finite extension of characteristic-zero fields,
of degree $n>1$. Its norm and trace are the determinant and trace
of multiplication on the $K$-vector space $L$. Define
\[
\begin{split}
 G_N&=\{(a,b)\in L^2:b=\Nm_{L/K}(a)\},\\
 G_T&=\{(a,b)\in L^2:b=\Tr_{L/K}(a)\}.
\end{split}
\]
These are graphs on one $L$-valued input. A norm polynomial in
separate $K$-valued basis coordinates is a different presentation:
its use inside $L$ requires restrictions on those coordinates.
Existential definitions with such restrictions are studied
in~\cite{FieldDefinitions}. Here a quantifier-free definition means a
finite Boolean combination of polynomial equalities with coefficients
in $L$ and variables ranging over $L$.

\begin{proposition}\label{prop:dense}
If $L/K$ is proper, the graphs $G_N$ and $G_T$ are Zariski dense
in $\A^2_L$. Neither has a quantifier-free definition in the
language of rings with coefficients in $L$. In particular neither
is the graph of a polynomial function in one variable over $L$.
\end{proposition}
\begin{proof}
Choose $\theta\in L\setminus K$. The set $K^2$ is Zariski dense
in $\A^2_L$, because a polynomial vanishing on $K^2$ is zero by
applying the one-variable root bound successively.

Consider first the trace graph. Its points include the images of
$K^2$ under the $L$-linear map
\[
    (a,b)\longmapsto
        \bigl(a+b\theta,\;na+b\Tr_{L/K}(\theta)\bigr).
\]
Its determinant is $\Tr_{L/K}(\theta)-n\theta\ne0$:
otherwise $\theta\in K$. Hence this map is an automorphism of
$\A^2_L$, proving density of $G_T$.

For the norm, restrict the polynomial norm form to the two
coordinates $a,b$ corresponding to $a+b\theta$ and consider
\[
        (a,b)\longmapsto
          \bigl(a+b\theta,\;\Nm_{L/K}(a+b\theta)\bigr).
\]
At $(1,0)$ its Jacobian determinant is again
$\Tr_{L/K}(\theta)-n\theta\ne0$, since the linear term in
$\Nm(1+b\theta)$ is $b\Tr(\theta)$. The map is dominant.
Its values on $K^2$ belong to $G_N$, so $G_N$ is dense as well.

A quantifier-free formula defines a constructible set. A dense
constructible subset of $\A^2_L$ contains a nonempty Zariski
open subset. Since $L$ is infinite, such an open subset has
infinitely many points in at least one vertical line; for instance,
choose a first coordinate at which a nonzero defining polynomial
of a principal open subset remains nonzero. A graph of a function
has exactly one point in each vertical line. This contradiction
rules out a quantifier-free definition. A polynomial graph is a
proper closed subset and is excluded in particular.
\end{proof}

\begin{thebibliography}{99}
\bibitem{FieldDefinitions}
\emph{Existential Definitions of Rings of Integers and Subfields},
project manuscript, 9 November 2026,
\nolinkurl{papers/number_fields/number_field_definability_20261109_v12.tex}.
\end{thebibliography}
\end{document}
